% Figure: a compound-lever bolt cutter. The long handles drive a short link
% that in turn drives the cutting jaw about its pivot, two levers in series
% multiplying the hand force twice. The two stages are labelled with their
% arm ratios; the product is the tool's mechanical advantage.
\begin{figure}[htbp]
\centering
\begin{tikzpicture}[>=latex, x=1.0cm, y=1.0cm]
  % stage 1: handle lever

  % handle pivot
  \coordinate (P1) at (0,0);
  % handle far end (hand)
  \coordinate (H) at (5.2,-0.5);
  % link connection point on handle (short arm)
  \coordinate (K) at (0.9,0.25);
  \draw[engblue, very thick] (P1) -- (H);
  \fill[engblack] (P1) circle (2pt);
  \node[font=\scriptsize, anchor=north east] at (P1) {pivot $O_1$};
  \draw[->, very thick, engorange] (5.0,0.6) -- (5.0,-0.3);
  \node[engorange, font=\scriptsize, anchor=west] at (5.1,0.2) {$P$ (hand)};
  \fill[engred] (K) circle (1.8pt);
  \draw[<->, enggray, thin] (0,-0.9) -- (0.9,-0.9);
  \node[enggray, font=\scriptsize, anchor=north] at (0.45,-0.9) {$a_1$};
  \draw[<->, enggray, thin] (0,-1.5) -- (5.2,-1.5);
  \node[enggray, font=\scriptsize, anchor=north] at (2.6,-1.5) {$b_1$};
  % link

  \coordinate (J) at (1.7,1.5); % jaw drive point
  \draw[engamber, very thick] (K) -- (J);
  \node[engamber, font=\scriptsize, anchor=west] at (1.15,0.9) {link};
  % stage 2: jaw lever

  \coordinate (P2) at (2.6,1.05); % jaw pivot
  \coordinate (C) at (0.9,2.4);   % cutting edge
  \draw[engblue, very thick] (P2) -- (C);
  \draw[engblue, very thick] (P2) -- (J);
  \fill[engblack] (P2) circle (2pt);
  \node[font=\scriptsize, anchor=west] at (P2) {pivot $O_2$};
  \fill[engred] (J) circle (1.8pt);
  % workpiece (bolt) at the cutting edge
  \draw[enggray, very thick] (0.6,2.55) -- (1.2,2.25);
  \draw[->, very thick, engred] (0.3,2.9) -- (0.75,2.5);
  \node[engred, font=\scriptsize, anchor=south] at (0.3,2.9) {$Q$ (cut)};
  \draw[<->, enggray, thin] (2.9,1.2) -- (2.05,1.6);
  \node[enggray, font=\scriptsize, anchor=south west] at (2.4,1.35) {$a_2$};
  \draw[<->, enggray, thin] (2.9,1.35) -- (1.15,2.55);
  \node[enggray, font=\scriptsize, anchor=west] at (2.1,2.1) {$b_2$};
\end{tikzpicture}
\caption[Compound-lever bolt cutter]{A bolt cutter's compound lever. The hand
force $P$ acts at the long arm $b_1$ of the handle, which pivots at $O_1$ and
drives a short link from its short arm $a_1$; the link in turn drives the
cutting jaw at the short arm $a_2$ of a second lever pivoted at $O_2$, whose
cutting edge sits at the short arm $b_2$. Each stage multiplies the force by
its arm ratio, and the two stages act in series, so the overall mechanical
advantage is the product $(b_1/a_1)\,(a_2/b_2)$ of the two ratios. Two modest
lever ratios multiplied give the large advantage that lets a hand force part a
hardened bolt, the compound-lever principle behind bolt cutters, pruning
loppers, and crimping tools.}
\label{fig:vol03ch02:bolt-cutter}
\end{figure}
